\subsection{Modular Arithmetic Continued}

\content{
        Let's review what we have learned so far.
}{% presentation
\newpage
}{% nopresentation
\vspace{4in}
}{% comments
\begin{enumerate}
\item Start with definition of modular arithmetic. 
\item Go over addition example, show that any number which has the same
remainder can be used in a numbers place
\item Go over multiplication example, same thing
\item Go over exponentiaion example, same thing.
\end{enumerate}
}

\content{
        \begin{theorem} (Modular Exponent Power Rule) when $n$ is a natural
        number, we have 
        $$ (a^b\mod{n})^c\mod{n} = (a^{bc}) \mod{n} = (a^c\mod{n})^b\mod{n}. $$
        \end{theorem}

        We also note, that when $p$ is a prime number, it is always possible to
        find a `generator', that is, a number $g$ such that all the powers of
        $g$ modulo $p$ give all the possible values from 1 to $p-1$. 

        \begin{example} Compute $3^n\mod{7}$ for all values of $n$.
        \end{example}
}{% presentation
\vspace{3in}
}{% nopresentation
\vspace{3in}
}{% comments
}

\content{
        We are now ready to describe the key exchange algorithm. 
        \begin{enumerate}
        \item Alice and Bob agree publically on values a prime $p$ and generator
        $g$.
        \item Alice picks some secret number $a$, while Bob picks some secret
        number $b$.
        \item Alice computes $A = g^a \mod{p}$ and sends it to Bob.
        \item Bob computes $B = g^b \mod{p}$ and sends it to Alice.
        \item Alice computes $B^a \mod{p}$, which is $(g^b \mod{p})^a \mod{p}$
        \item Bob computes $A^b \mod{p}$, which is $(g^a \mod{p})^b \mod{p}$
        \end{enumerate}
}{% presentation
}{% nopresentation
}{% comments
}

\content{
        \begin{example} Suppose that Alice and Bob publicly choose the prime
        $17$ and generator $3$. What does the exchange look like, and where is
        their secret key?
        \end{example}
}{% presentation
\vspace{2in}
}{% nopresentation
\vspace{2in}
}{% comments
}


\content{
        \begin{example} Go through the key exchange algorithm with public
        values $p=17$ and generator $g=3$.
        \end{example}
}{% presentation
\vspace{3in}
}{% nopresentation
\vspace{3in}
}{% comments
}

\content{
        \begin{example} Your turn! What does the key exchange algorithm look
        like with $p=19$ and $g=13$?
        \end{example}
}{% presentation
\vspace{3in}
}{% nopresentation
\vspace{3in}
}{% comments
}
